% Kalkylark reference card: one A4 landscape page, three columns.
% Build: make docs   (tectonic). Licensed CC BY 4.0.
\documentclass[10pt]{article}
\PassOptionsToPackage{landscape}{geometry}
\usepackage{kalkylark}
\usepackage{multicol}
\newgeometry{landscape, margin=11mm}
\pagestyle{empty}
\setlength{\columnsep}{9mm}
\setlength{\columnseprule}{0.3pt}
\renewcommand{\columnseprulecolor}{\color{phdim}}
\setlength{\parskip}{2pt}

% compact sections and two-column key tables
\newcommand{\rc}[1]{\par\vspace{5pt}{\sffamily\bfseries\color{ink}#1}\par\vspace{1pt}}
\newenvironment{keys}[1][22mm]
  {\par\small\tabularx{\linewidth}{@{}>{\raggedright\arraybackslash}p{#1}>{\raggedright\arraybackslash}X@{}}}
  {\endtabularx\par}
\newcommand{\kt}[1]{\texttt{#1}}

\begin{document}
\raggedright
\raggedcolumns
\begin{multicols*}{3}

{\sffamily\bfseries\Huge\color{ink}Kalkylark}\par
{\sffamily\large\color{ink}Reference card · version 0.1}\par
\vspace{4pt}
{\small A VisiCalc-style spreadsheet for VT100 terminals.
63 columns (\kt{A}--\kt{BK}) × 254 rows.}

\rc{Starting}
\begin{keys}[30mm]
\kt{kalkylark} & empty sheet \\
\kt{kalkylark f.kalk} & open or create \kt{f.kalk} \\
\kt{-p green} & phosphor: \kt{plain}, \kt{green}, \kt{amber}, \kt{white} \\
\kt{-c 256} & colours: \kt{16}, \kt{256}, \kt{true} \\
\kt{-V} & version \\
\end{keys}
Environment: \kt{KALKYLARK\_PHOSPHOR}, \kt{NO\_COLOR}.

\rc{The screen}
\begin{center}
\setlength{\fboxsep}{4pt}
\colorbox{phbg}{\begin{minipage}{0.92\linewidth}\ttfamily\scriptsize
\textcolor{phbright}{B3 /F\$ (V) @SUM(B1...B2)\hfill}\textcolor{phdim}{N KALKYLARK}\\
\textcolor{phnorm}{prompt and messages}\\
\textcolor{phbright}{entry being typed}\textcolor{phbright}{\rule{0.5em}{0.8em}}\\
{\setlength{\fboxsep}{0pt}\colorbox{phdim}{\textcolor{phbg}{\strut\ \ \ \ \ \ \ \ \ A\ \ \ \ \ \ \ \ B\ \ \ \ \ \ \ \ C\ \ \ }}}
\end{minipage}}
\end{center}
\vspace{-4pt}
{\small Status line: cell, format, \textsf{(L)}abel or \textsf{(V)}alue,
contents; at the right \textsf{CIRC} if there is a circular reference,
then the recalculation order \textsf{C}/\textsf{R}/\textsf{N}.}

\rc{Moving}
\begin{keys}
\up\ \down\ \lft\ \rgt & one cell \\
\kbd{PgUp} \kbd{PgDn} & one screen \\
\kbd{Home} & to \kt{A1} \\
\kt{>}\emph{cell}\enter & go to, e.g. \kt{>BK254} \\
\kbd{Ctrl-L} & redraw \\
\end{keys}

\rc{Entering}
The \textbf{first key} decides:
\begin{keys}
\kt{0-9 + - . ( @} & starts a \textbf{value} \\
anything else & starts a \textbf{label} \\
\kt{"}\emph{text} & forces a label: \kt{"2026} \\
\end{keys}
\begin{keys}
\enter & store \\
arrow & store and move \\
\kbd{Backspace} & delete a character \\
\esc & cancel \\
\end{keys}

\rc{Pointing}
After \kt{+ - * / \^{} ( ,} or \kt{.} the arrows move a \emph{pointer} and
insert its coordinate. Typing \kt{.} while pointing writes \kt{...} and
starts the range's second corner:\par
\kt{@SUM(} \up\up\ \kt{.} \down\down\ \kt{)} \(\to\) \kt{@SUM(B1...B3)}

\columnbreak

\rc{Formulas}
\textbf{Strictly left to right, no precedence:} \kt{2+3*4} = 20,
\kt{2+(3*4)} = 14.
\begin{keys}[24mm]
\kt{+ - * /} & arithmetic; \kt{/0} gives \textsf{ERROR} \\
\kt{\^{}} & power \\
\kt{-x} & negation of the operand \\
\kt{1.5E3} & scientific notation \\
\kt{A1} & reference; blank and label = 0 \\
\kt{A1...C9} & range; also \kt{A1.C9} \kt{A1..C9} \kt{A1:C9} \\
\end{keys}
A formula starting with a cell needs \kt{+}: \kt{+A1*2}.\par
Careful: \kt{1-(1+r)\^{}n} means \((1-(1+r))^n\).
Write \kt{1-((1+r)\^{}n)}.

\rc{Functions}
\begin{keys}[27mm]
\kt{@SUM(list)} & sum \\
\kt{@MIN(list)} \kt{@MAX} & smallest, largest \\
\kt{@COUNT(list)} & number of values \\
\kt{@AVERAGE(list)} & mean \\
\kt{@NPV(r,range)} & \(\sum v_i/(1+r)^i\), \(i=1\dots\) \\
\kt{@LOOKUP(x,range)} & last entry \(\le x\) in an ascending row or column;
  result from the cell to its right (or below) \\
\kt{@ABS @INT} & absolute value, integer part \\
\kt{@SQRT @EXP} & square root, \(e^x\) \\
\kt{@LN @LOG10} & logarithms \\
\kt{@SIN @COS @TAN} & radians \\
\kt{@ASIN @ACOS @ATAN} & inverses \\
\kt{@PI} & \(\pi\) \\
\kt{@NA} & ``not available'' \\
\kt{@ERROR} & an error value \\
\end{keys}
A \emph{list} mixes values and ranges: \kt{@SUM(A1...A9,C1,100)}.
In ranges, blanks and labels are skipped.
\textsf{NA} and \textsf{ERROR} spread through formulas; \textsf{ERROR}
wins.

\rc{Formats}
\begin{keys}[8mm]
\kt{G} & general: as many digits as fit \\
\kt{I} & integer, rounded \\
\kt{\$} & two decimals \\
\kt{L} & left-aligned \\
\kt{R} & right-aligned (labels too) \\
\kt{*} & bar graph: one \kt{*} per unit \\
\kt{D} & default: use the global format \\
\end{keys}
\kt{>>>>} means the number does not fit: widen with \kt{/GC}.

\columnbreak

\rc{Commands}
Type \kt{/}, then letters. \esc{} backs out.
\begin{keys}[19mm]
\kt{/B}\enter & blank the cell \\
\kt{/C Y} & clear the sheet \\
\kt{/F}\emph{f} & cell format (\kt{G I \$ L R * D}) \\
\kt{/GF}\emph{f} & global format \\
\kt{/GC}\emph{n}\enter & column width 3--40, all columns \\
\kt{/GOC} & recalc by column (default) \\
\kt{/GOR} & recalc by row \\
\kt{/GON} & recalc in natural order \\
\kt{/SS} & save \\
\kt{/SL} & load \\
\kt{/SQY} & quit \\
\end{keys}
In name and number prompts the suggested value is accepted with \enter,
replaced by typing, and edited with \kbd{Backspace}.

\rc{Recalculation}
\begin{keys}
\kt{!} & recalculate now \\
\end{keys}
\textbf{C} and \textbf{R} (VisiCalc): one pass in column or row order.
A formula that refers forward may show an old value. Press \kt{!} again.\par
\textbf{N} (Lotus 1-2-3): each formula after the ones it reads, so one
recalculation is always enough. \textsf{CIRC} marks a circular reference;
its cells move one step per \kt{!}.

\rc{Files}
Plain text, one record per line:
\begin{lstlisting}[basicstyle=\ttfamily\footnotesize, frame=none,
  xleftmargin=0pt, framexleftmargin=0pt]
# kalkylark 1
/GO N
/GC 9
/GF G
>A1:"Rent
>B1:/F$:9500
>B3:@SUM(B1...B2)
\end{lstlisting}
Saves are atomic (temporary file, then rename). Kalkylark asks before
overwriting, and before quitting or loading over unsaved changes.

\vfill
{\scriptsize\color{ink}Kalkylark is free software (BSD 2-clause) by Set
Lonnert and the Kalkylark authors. This card: CC BY 4.0. Not affiliated
with Lotus, IBM, Software Arts or VisiCorp. VisiCalc is a trademark of its
respective owner.\par}

\end{multicols*}
\end{document}
